\IfFileExists{stacks-project.cls}{%
\documentclass{stacks-project}
}{%
\documentclass{amsart}
}

% For dealing with references we use the comment environment
\usepackage{verbatim}
\newenvironment{reference}{\comment}{\endcomment}
%\newenvironment{reference}{}{}
\newenvironment{slogan}{\comment}{\endcomment}
\newenvironment{history}{\comment}{\endcomment}

% For commutative diagrams we use Xy-pic
\usepackage[all]{xy}

% We use 2cell for 2-commutative diagrams.
\xyoption{2cell}
\UseAllTwocells

% We use multicol for the list of chapters between chapters
\usepackage{multicol}

% This is generally recommended for better output
\usepackage{lmodern}
\usepackage[T1]{fontenc}

% For cross-file-references
\usepackage{xr-hyper}

% Package for hypertext links:
\usepackage{hyperref}

% For any local file, say "hello.tex" you want to link to please
% use \externaldocument[hello-]{hello}
\externaldocument[introduction-]{introduction}
\externaldocument[conventions-]{conventions}
\externaldocument[sets-]{sets}
\externaldocument[categories-]{categories}
\externaldocument[topology-]{topology}
\externaldocument[sheaves-]{sheaves}
\externaldocument[sites-]{sites}
\externaldocument[stacks-]{stacks}
\externaldocument[fields-]{fields}
\externaldocument[algebra-]{algebra}
\externaldocument[brauer-]{brauer}
\externaldocument[homology-]{homology}
\externaldocument[derived-]{derived}
\externaldocument[simplicial-]{simplicial}
\externaldocument[more-algebra-]{more-algebra}
\externaldocument[smoothing-]{smoothing}
\externaldocument[modules-]{modules}
\externaldocument[sites-modules-]{sites-modules}
\externaldocument[injectives-]{injectives}
\externaldocument[cohomology-]{cohomology}
\externaldocument[sites-cohomology-]{sites-cohomology}
\externaldocument[dga-]{dga}
\externaldocument[dpa-]{dpa}
\externaldocument[sdga-]{sdga}
\externaldocument[hypercovering-]{hypercovering}
\externaldocument[schemes-]{schemes}
\externaldocument[constructions-]{constructions}
\externaldocument[properties-]{properties}
\externaldocument[morphisms-]{morphisms}
\externaldocument[coherent-]{coherent}
\externaldocument[divisors-]{divisors}
\externaldocument[limits-]{limits}
\externaldocument[varieties-]{varieties}
\externaldocument[topologies-]{topologies}
\externaldocument[descent-]{descent}
\externaldocument[perfect-]{perfect}
\externaldocument[more-morphisms-]{more-morphisms}
\externaldocument[flat-]{flat}
\externaldocument[groupoids-]{groupoids}
\externaldocument[more-groupoids-]{more-groupoids}
\externaldocument[etale-]{etale}
\externaldocument[chow-]{chow}
\externaldocument[intersection-]{intersection}
\externaldocument[pic-]{pic}
\externaldocument[weil-]{weil}
\externaldocument[adequate-]{adequate}
\externaldocument[dualizing-]{dualizing}
\externaldocument[duality-]{duality}
\externaldocument[discriminant-]{discriminant}
\externaldocument[derham-]{derham}
\externaldocument[local-cohomology-]{local-cohomology}
\externaldocument[algebraization-]{algebraization}
\externaldocument[curves-]{curves}
\externaldocument[resolve-]{resolve}
\externaldocument[models-]{models}
\externaldocument[functors-]{functors}
\externaldocument[equiv-]{equiv}
\externaldocument[pione-]{pione}
\externaldocument[etale-cohomology-]{etale-cohomology}
\externaldocument[proetale-]{proetale}
\externaldocument[relative-cycles-]{relative-cycles}
\externaldocument[more-etale-]{more-etale}
\externaldocument[trace-]{trace}
\externaldocument[crystalline-]{crystalline}
\externaldocument[spaces-]{spaces}
\externaldocument[spaces-properties-]{spaces-properties}
\externaldocument[spaces-morphisms-]{spaces-morphisms}
\externaldocument[decent-spaces-]{decent-spaces}
\externaldocument[spaces-cohomology-]{spaces-cohomology}
\externaldocument[spaces-limits-]{spaces-limits}
\externaldocument[spaces-divisors-]{spaces-divisors}
\externaldocument[spaces-over-fields-]{spaces-over-fields}
\externaldocument[spaces-topologies-]{spaces-topologies}
\externaldocument[spaces-descent-]{spaces-descent}
\externaldocument[spaces-perfect-]{spaces-perfect}
\externaldocument[spaces-more-morphisms-]{spaces-more-morphisms}
\externaldocument[spaces-flat-]{spaces-flat}
\externaldocument[spaces-groupoids-]{spaces-groupoids}
\externaldocument[spaces-more-groupoids-]{spaces-more-groupoids}
\externaldocument[bootstrap-]{bootstrap}
\externaldocument[spaces-pushouts-]{spaces-pushouts}
\externaldocument[spaces-chow-]{spaces-chow}
\externaldocument[groupoids-quotients-]{groupoids-quotients}
\externaldocument[spaces-more-cohomology-]{spaces-more-cohomology}
\externaldocument[spaces-simplicial-]{spaces-simplicial}
\externaldocument[spaces-duality-]{spaces-duality}
\externaldocument[formal-spaces-]{formal-spaces}
\externaldocument[restricted-]{restricted}
\externaldocument[spaces-resolve-]{spaces-resolve}
\externaldocument[formal-defos-]{formal-defos}
\externaldocument[defos-]{defos}
\externaldocument[cotangent-]{cotangent}
\externaldocument[examples-defos-]{examples-defos}
\externaldocument[algebraic-]{algebraic}
\externaldocument[examples-stacks-]{examples-stacks}
\externaldocument[stacks-sheaves-]{stacks-sheaves}
\externaldocument[criteria-]{criteria}
\externaldocument[artin-]{artin}
\externaldocument[quot-]{quot}
\externaldocument[stacks-properties-]{stacks-properties}
\externaldocument[stacks-morphisms-]{stacks-morphisms}
\externaldocument[stacks-limits-]{stacks-limits}
\externaldocument[stacks-cohomology-]{stacks-cohomology}
\externaldocument[stacks-perfect-]{stacks-perfect}
\externaldocument[stacks-introduction-]{stacks-introduction}
\externaldocument[stacks-more-morphisms-]{stacks-more-morphisms}
\externaldocument[stacks-geometry-]{stacks-geometry}
\externaldocument[moduli-]{moduli}
\externaldocument[moduli-curves-]{moduli-curves}
\externaldocument[examples-]{examples}
\externaldocument[exercises-]{exercises}
\externaldocument[guide-]{guide}
\externaldocument[desirables-]{desirables}
\externaldocument[coding-]{coding}
\externaldocument[obsolete-]{obsolete}
\externaldocument[fdl-]{fdl}
\externaldocument[index-]{index}

% Theorem environments.
%
\theoremstyle{plain}
\newtheorem{theorem}[subsection]{Theorem}
\newtheorem{proposition}[subsection]{Proposition}
\newtheorem{lemma}[subsection]{Lemma}

\theoremstyle{definition}
\newtheorem{definition}[subsection]{Definition}
\newtheorem{example}[subsection]{Example}
\newtheorem{exercise}[subsection]{Exercise}
\newtheorem{situation}[subsection]{Situation}

\theoremstyle{remark}
\newtheorem{remark}[subsection]{Remark}
\newtheorem{remarks}[subsection]{Remarks}

\numberwithin{equation}{subsection}

% Macros
%
\def\lim{\mathop{\mathrm{lim}}\nolimits}
\def\colim{\mathop{\mathrm{colim}}\nolimits}
\def\Spec{\mathop{\mathrm{Spec}}}
\def\Hom{\mathop{\mathrm{Hom}}\nolimits}
\def\Ext{\mathop{\mathrm{Ext}}\nolimits}
\def\SheafHom{\mathop{\mathcal{H}\!\mathit{om}}\nolimits}
\def\SheafExt{\mathop{\mathcal{E}\!\mathit{xt}}\nolimits}
\def\Sch{\mathit{Sch}}
\def\Mor{\mathop{\mathrm{Mor}}\nolimits}
\def\Ob{\mathop{\mathrm{Ob}}\nolimits}
\def\Sh{\mathop{\mathit{Sh}}\nolimits}
\def\NL{\mathop{N\!L}\nolimits}
\def\CH{\mathop{\mathrm{CH}}\nolimits}
\def\proetale{{pro\text{-}\acute{e}tale}}
\def\etale{{\acute{e}tale}}
\def\QCoh{\mathit{QCoh}}
\def\Ker{\mathop{\mathrm{Ker}}}
\def\Im{\mathop{\mathrm{Im}}}
\def\Coker{\mathop{\mathrm{Coker}}}
\def\Coim{\mathop{\mathrm{Coim}}}

% Boxtimes
%
\DeclareMathSymbol{\boxtimes}{\mathbin}{AMSa}{"02}

%
% Macros for moduli stacks/spaces
%
\def\QCohstack{\mathcal{QC}\!\mathit{oh}}
\def\Cohstack{\mathcal{C}\!\mathit{oh}}
\def\Spacesstack{\mathcal{S}\!\mathit{paces}}
\def\Quotfunctor{\mathrm{Quot}}
\def\Hilbfunctor{\mathrm{Hilb}}
\def\Curvesstack{\mathcal{C}\!\mathit{urves}}
\def\Polarizedstack{\mathcal{P}\!\mathit{olarized}}
\def\Complexesstack{\mathcal{C}\!\mathit{omplexes}}
% \Pic is the operator that assigns to X its picard group, usage \Pic(X)
% \Picardstack_{X/B} denotes the Picard stack of X over B
% \Picardfunctor_{X/B} denotes the Picard functor of X over B
\def\Pic{\mathop{\mathrm{Pic}}\nolimits}
\def\Picardstack{\mathcal{P}\!\mathit{ic}}
\def\Picardfunctor{\mathrm{Pic}}
\def\Deformationcategory{\mathcal{D}\!\mathit{ef}}
\usepackage{cleveref}
\crefname{originalreference}{original reference}{original references}
\Crefname{originalreference}{Original reference}{Original references}
\makeatletter
\newlabel{algebra-lemma-NAK}{{00DV}{source}{}{sourcepointer-1}{}}
\newlabel{algebra-lemma-NAK@cref}{{[originalreference][0][]00DV}{[1][0][]source}{}{}{}}
\makeatother
\crefname{originalequation}{original equation}{original equations}
\Crefname{originalequation}{Original equation}{Original equations}
\makeatletter
\newlabel{equation-gkn}{{089S}{source}{}{sourcepointer-2}{}}
\newlabel{equation-gkn@cref}{{[originalequation][0][]089S}{[1][0][]source}{}{}{}}
\makeatother
\crefname{originalreference}{original reference}{original references}
\Crefname{originalreference}{Original reference}{Original references}
\makeatletter
\newlabel{modules-lemma-exactness-pushforward-pullback}{{01AJ}{source}{}{sourcepointer-3}{}}
\newlabel{modules-lemma-exactness-pushforward-pullback@cref}{{[originalreference][0][]01AJ}{[1][0][]source}{}{}{}}
\makeatother
\crefname{originalreference}{original reference}{original references}
\Crefname{originalreference}{Original reference}{Original references}
\makeatletter
\newlabel{modules-lemma-finite-type-surjective-on-stalk}{{01B8}{source}{}{sourcepointer-4}{}}
\newlabel{modules-lemma-finite-type-surjective-on-stalk@cref}{{[originalreference][0][]01B8}{[1][0][]source}{}{}{}}
\makeatother
\crefname{originalreference}{original reference}{original references}
\Crefname{originalreference}{Original reference}{Original references}
\makeatletter
\newlabel{modules-lemma-map-finite-locally-free}{{089Q}{source}{}{sourcepointer-5}{}}
\newlabel{modules-lemma-map-finite-locally-free@cref}{{[originalreference][0][]089Q}{[1][0][]source}{}{}{}}
\makeatother
\crefname{originalreference}{original reference}{original references}
\Crefname{originalreference}{Original reference}{Original references}
\makeatletter
\newlabel{schemes-example-global-sections}{{01JH}{source}{}{sourcepointer-6}{}}
\newlabel{schemes-example-global-sections@cref}{{[originalreference][0][]01JH}{[1][0][]source}{}{}{}}
\makeatother
\crefname{originalreference}{original reference}{original references}
\Crefname{originalreference}{Original reference}{Original references}
\makeatletter
\newlabel{schemes-lemma-glue-functors}{{01JJ}{source}{}{sourcepointer-7}{}}
\newlabel{schemes-lemma-glue-functors@cref}{{[originalreference][0][]01JJ}{[1][0][]source}{}{}{}}
\makeatother
\crefname{originalreference}{original reference}{original references}
\Crefname{originalreference}{Original reference}{Original references}
\makeatletter
\newlabel{sheaves-lemma-stalk-pullback-modules}{{0098}{source}{}{sourcepointer-8}{}}
\newlabel{sheaves-lemma-stalk-pullback-modules@cref}{{[originalreference][0][]0098}{[1][0][]source}{}{}{}}
\makeatother
\crefname{originalreference}{original reference}{original references}
\Crefname{originalreference}{Original reference}{Original references}
\makeatletter
\newlabel{sheaves-section-glueing-sheaves}{{00AK}{source}{}{sourcepointer-9}{}}
\newlabel{sheaves-section-glueing-sheaves@cref}{{[originalreference][0][]00AK}{[1][0][]source}{}{}{}}
\makeatother
\title{The Grassmannian functor G(k, n) of quotients is representable}
\author{The Stacks Project Authors}
\date{stacks-project:constructions.tex:local-snapshot-20261012; GNU FDL 1.2 or later}
\begin{document}
\maketitle
\noindent\textit{Editorial note.} Opening sheaf-property sentence carries 'some details omitted' (glueing of quotient data via Sheaves, Section glueing-sheaves); the four main steps are fully displayed.\par\medskip
\noindent\textit{Editorial note.} Steps the author leaves to the reader are flagged here for the reviewer: source line(s) 5221 carry an "omitted"/"follows easily"/"immediate from the definitions" phrase. Those sentences are retained verbatim; the displayed argument around them stays inside the counted span.\par\medskip
\noindent\textit{Allowed standard prerequisites (editorial).}\begin{itemize}
\item Glue-functors criterion. schemes-lemma-glue-functors (cross-chapter, named). Provides the four-step representability template used for the whole proof.
\item Right exactness of pullback. modules-lemma-exactness-pushforward-pullback (cross-chapter, named). Makes each chart subfunctor F\_i stable under pullback.
\item Surjectivity on stalks is open. modules-lemma-finite-type-surjective-on-stalk (cross-chapter, named). Opens up the chart locus U\_i inside S.
\item Nakayama. algebra-lemma-NAK (cross-chapter, named). Reduces chart membership to surjectivity on the residue fibre.
\end{itemize}\medskip
\ifdefined\hypertarget\hypertarget{sourcepointer-1}{}\fi
\noindent\textit{Source reference.} Original reference 00DV: focused-seeds/geometry/S9-cons-gkn-representable/ref-tags.txt, source lines 5--5. Navigation pointer only: the counted proof cites this statement by \textbackslash{}ref as a standing input; the deduction the counted proof has to make is carried out inside the counted span, and no part of that deduction is delegated to this pointer.\par\medskip
\ifdefined\hypertarget\hypertarget{sourcepointer-2}{}\fi
\noindent\textit{Source reference.} Original equation 089S: sources/stacks-constructions/pinned-original-source.tex, source lines 5180--5180. Navigation pointer only: the counted proof cites this statement by \textbackslash{}ref as a standing input; the deduction the counted proof has to make is carried out inside the counted span, and no part of that deduction is delegated to this pointer.\par\medskip
\ifdefined\hypertarget\hypertarget{sourcepointer-3}{}\fi
\noindent\textit{Source reference.} Original reference 01AJ: focused-seeds/geometry/S9-cons-gkn-representable/ref-tags.txt, source lines 6--6. Navigation pointer only: the counted proof cites this statement by \textbackslash{}ref as a standing input; the deduction the counted proof has to make is carried out inside the counted span, and no part of that deduction is delegated to this pointer.\par\medskip
\ifdefined\hypertarget\hypertarget{sourcepointer-4}{}\fi
\noindent\textit{Source reference.} Original reference 01B8: focused-seeds/geometry/S9-cons-gkn-representable/ref-tags.txt, source lines 7--7. Navigation pointer only: the counted proof cites this statement by \textbackslash{}ref as a standing input; the deduction the counted proof has to make is carried out inside the counted span, and no part of that deduction is delegated to this pointer.\par\medskip
\ifdefined\hypertarget\hypertarget{sourcepointer-5}{}\fi
\noindent\textit{Source reference.} Original reference 089Q: focused-seeds/geometry/S9-cons-gkn-representable/ref-tags.txt, source lines 8--8. Navigation pointer only: the counted proof cites this statement by \textbackslash{}ref as a standing input; the deduction the counted proof has to make is carried out inside the counted span, and no part of that deduction is delegated to this pointer.\par\medskip
\ifdefined\hypertarget\hypertarget{sourcepointer-6}{}\fi
\noindent\textit{Source reference.} Original reference 01JH: focused-seeds/geometry/S9-cons-gkn-representable/ref-tags.txt, source lines 9--9. Navigation pointer only: the counted proof cites this statement by \textbackslash{}ref as a standing input; the deduction the counted proof has to make is carried out inside the counted span, and no part of that deduction is delegated to this pointer.\par\medskip
\ifdefined\hypertarget\hypertarget{sourcepointer-7}{}\fi
\noindent\textit{Source reference.} Original reference 01JJ: focused-seeds/geometry/S9-cons-gkn-representable/ref-tags.txt, source lines 10--10. Navigation pointer only: the counted proof cites this statement by \textbackslash{}ref as a standing input; the deduction the counted proof has to make is carried out inside the counted span, and no part of that deduction is delegated to this pointer.\par\medskip
\ifdefined\hypertarget\hypertarget{sourcepointer-8}{}\fi
\noindent\textit{Source reference.} Original reference 0098: focused-seeds/geometry/S9-cons-gkn-representable/ref-tags.txt, source lines 11--11. Navigation pointer only: the counted proof cites this statement by \textbackslash{}ref as a standing input; the deduction the counted proof has to make is carried out inside the counted span, and no part of that deduction is delegated to this pointer.\par\medskip
\ifdefined\hypertarget\hypertarget{sourcepointer-9}{}\fi
\noindent\textit{Source reference.} Original reference 00AK: focused-seeds/geometry/S9-cons-gkn-representable/ref-tags.txt, source lines 12--12. Navigation pointer only: the counted proof cites this statement by \textbackslash{}ref as a standing input; the deduction the counted proof has to make is carried out inside the counted span, and no part of that deduction is delegated to this pointer.\par\medskip
\begin{lemma}
\label{lemma-gkn-representable}
Let $0 < k < n$.
The functor $G(k, n)$ of (\ref{equation-gkn}) is representable by a scheme.
\end{lemma}

\begin{proof}
Set $F = G(k, n)$. To prove the lemma we will use the criterion of
Schemes, Lemma \ref{schemes-lemma-glue-functors}.
The reason $F$ satisfies the sheaf property for the
Zariski topology is that we can glue sheaves, see Sheaves,
Section \ref{sheaves-section-glueing-sheaves} (some details omitted).

\medskip\noindent
The family of subfunctors $F_i$.
Let $I$ be the set of subsets of $\{1, \ldots, n\}$ of cardinality $n - k$.
Given a scheme $S$ and $j \in \{1, \ldots, n\}$ we denote $e_j$
the global section
$$
e_j = (0, \ldots, 0, 1, 0, \ldots, 0)\quad(1\text{ in }j\text{th spot})
$$
of $\mathcal{O}_S^{\oplus n}$. Of course these sections freely generate
$\mathcal{O}_S^{\oplus n}$. Similarly, for $j \in \{1, \ldots, n - k\}$
we denote $f_j$ the global section of $\mathcal{O}_S^{\oplus n - k}$
which is zero in all summands except the $j$th where we put a $1$.
For $i \in I$ we let
$$
s_i : \mathcal{O}_S^{\oplus n - k} \longrightarrow \mathcal{O}_S^{\oplus n}
$$
which is the direct sum of the coprojections
$\mathcal{O}_S \to \mathcal{O}_S^{\oplus n}$ corresponding to elements of $I$.
More precisely, if $i = \{i_1, \ldots, i_{n - k}\}$ with
$i_1 < i_2 < \ldots < i_{n - k}$
then $s_i$ maps $f_j$ to $e_{i_j}$ for $j \in \{1, \ldots, n - k\}$.
With this notation we can set
$$
F_i(S) = \{q : \mathcal{O}_S^{\oplus n} \to \mathcal{Q} \in F(S) \mid
q \circ s_i \text{ is surjective}\}
\subset F(S)
$$
Given a morphism $f : T \to S$ of schemes the pullback $f^*s_i$
is the corresponding map over $T$. Since $f^*$ is right exact
(Modules, Lemma \ref{modules-lemma-exactness-pushforward-pullback})
we conclude that $F_i$ is a subfunctor of $F$.

\medskip\noindent
Representability of $F_i$. To prove this we may assume (after renumbering)
that $i = \{1, \ldots, n - k\}$. This means $s_i$ is the inclusion of
the first $n - k$ summands. Observe that if $q \circ s_i$ is surjective,
then $q \circ s_i$ is an isomorphism as a surjective map between finite
locally free modules of the same rank
(Modules, Lemma \ref{modules-lemma-map-finite-locally-free}).
Thus if $q : \mathcal{O}_S^{\oplus n} \to \mathcal{Q}$ is an element of
$F_i(S)$, then we can use $q \circ s_i$ to identify $\mathcal{Q}$ with
$\mathcal{O}_S^{\oplus n - k}$. After doing so we obtain
$$
q : \mathcal{O}_S^{\oplus n} \longrightarrow \mathcal{O}_S^{\oplus n - k}
$$
mapping $e_j$ to $f_j$ (notation as above) for $j = 1, \ldots, n - k$.
To determine $q$ completely we have to fix the images
$q(e_{n - k + 1}), \ldots, q(e_n)$ in
$\Gamma(S, \mathcal{O}_S^{\oplus n - k})$.
It follows that $F_i$ is isomorphic to the functor
$$
S \longmapsto
\prod\nolimits_{j = n - k + 1, \ldots, n}
\Gamma(S,  \mathcal{O}_S^{\oplus n - k})
$$
This functor is isomorphic to the $k(n - k)$-fold self product of the functor
$S \mapsto \Gamma(S, \mathcal{O}_S)$. By
Schemes, Example \ref{schemes-example-global-sections}
the latter is representable by $\mathbf{A}^1_\mathbf{Z}$. It follows $F_i$
is representable by $\mathbf{A}^{k(n - k)}_\mathbf{Z}$ since fibred product
over $\Spec(\mathbf{Z})$ is the product in the category of schemes.

\medskip\noindent
The inclusion $F_i \subset F$ is representable by open immersions.
Let $S$ be a scheme and let
$q : \mathcal{O}_S^{\oplus n} \to \mathcal{Q}$ be an element of
$F(S)$. By
Modules, Lemma \ref{modules-lemma-finite-type-surjective-on-stalk}.
the set $U_i = \{s \in S \mid (q \circ s_i)_s\text{ surjective}\}$
is open in $S$. Since $\mathcal{O}_{S, s}$ is a local ring
and $\mathcal{Q}_s$ a finite $\mathcal{O}_{S, s}$-module
by Nakayama's lemma (Algebra, Lemma \ref{algebra-lemma-NAK}) we have
$$
s \in U_i \Leftrightarrow
\left(
\text{the map }
\kappa(s)^{\oplus n - k} \to \mathcal{Q}_s/\mathfrak m_s\mathcal{Q}_s
\text{ induced by }
(q \circ s_i)_s
\text{ is surjective}
\right)
$$
Let $f : T \to S$ be a morphism of schemes and let $t \in T$ be a point
mapping to $s \in S$. We have
$(f^*\mathcal{Q})_t =
\mathcal{Q}_s \otimes_{\mathcal{O}_{S, s}} \mathcal{O}_{T, t}$
(Sheaves, Lemma \ref{sheaves-lemma-stalk-pullback-modules})
and so on. Thus the map
$$
\kappa(t)^{\oplus n - k} \to (f^*\mathcal{Q})_t/\mathfrak m_t(f^*\mathcal{Q})_t
$$
induced by $(f^*q \circ f^*s_i)_t$ is the base change of the map
$\kappa(s)^{\oplus n - k} \to \mathcal{Q}_s/\mathfrak m_s\mathcal{Q}_s$
above by the field extension $\kappa(t)/\kappa(s)$. It follows
that $s \in U_i$ if and only if $t$ is in the corresponding open
for $f^*q$. In particular $T \to S$ factors through $U_i$ if
and only if $f^*q \in F_i(T)$ as desired.

\medskip\noindent
The collection $F_i$, $i \in I$ covers $F$. Let
$q : \mathcal{O}_S^{\oplus n} \to \mathcal{Q}$ be an element of
$F(S)$. We have to show that for every point $s$ of $S$ there exists
an $i \in I$ such that $s_i$ is surjective in a neighbourhood of $s$.
Thus we have to show that one of the compositions
$$
\kappa(s)^{\oplus n - k} \xrightarrow{s_i}
\kappa(s)^{\oplus n} \rightarrow
\mathcal{Q}_s/\mathfrak m_s\mathcal{Q}_s
$$
is surjective (see previous paragraph). As
$\mathcal{Q}_s/\mathfrak m_s\mathcal{Q}_s$ is a vector space of
dimension $n - k$ this follows from the theory of vector spaces.
\end{proof}

\end{document}
